% !PRINTED 421
\weberpage{421}

Ist nun $s$ eine beliebige zu $n$ theilerfremde Zahl, so können wir $h$ so bestimmen, dass
\[
h \equiv s \ (\mathrm{mod}\ n), \quad h \equiv 1 \ (\mathrm{mod}\ a)
\]
wird, und demnach folgt aus (9)
\[
\varphi(r^s, \alpha) = 0,
\]
d. h. die Gleichung (7) ist für \gesperrt{alle} Wurzeln von $X_n$ befriedigt. Es ist also $\varphi(x,\alpha)$ durch $X_n$ theilbar und muss also, da andererseits $X_n$ durch $\varphi(x,\alpha)$ theilbar vorausgesetzt war, bis auf einen numerischen Factor mit $X_n$ identisch sein.

\begin{center}
§. 135.

\gesperrt{Die Discriminante der Kreistheilungsgleichung.}
\end{center}

Die Einheitswurzeln vom Grade $n$ lassen sich in transcendenter Form darstellen durch
\weq{1}{r = e^{\frac{2\pi i k}{n}} = \cos\frac{2\pi k}{n} + i \sin\frac{2\pi k}{n}.}
Diese Einheitswurzel ist primitiv oder nicht primitiv, je nachdem $k$ theilerfremd zu $n$ ist oder nicht. Die einfachste unter den primitiven erhält man für $k = 1$, nämlich
\weq{2}{r_0 = e^{\frac{2\pi i}{n}} = \cos\frac{2\pi}{n} + i \sin\frac{2\pi}{n} = a + bi.}
Der Winkel $2\pi : n$ ist der $n^{\text{te}}$ Theil der ganzen Kreisperipherie. Die verschiedenen Theilpunkte, die man erhält, wenn man diesen Winkel von einem beliebigen Anfangspunkte an auf der Peripherie aufträgt, bestimmen das dem Kreise eingeschriebene reguläre $n$-Eck.

% Registered Bd. I scan, PDF page 447 / printed p. 421, rendered at 600 dpi.
% Exact source crop x=100,y=3150,w=810,h=550; SHA-256 27F4BCA9D93A7A0A2DED6313E3840AD29B999168A364F0FDE053D0091B6B361B.
\weberfigureimage{weber-b1-fig28.png}{0.28}

Die Theilpunkte lassen sich construiren, wenn man $r$ und damit $a$ und $b$ (oder auch nur eine dieser beiden Grössen) kennt. Insbesondere ergiebt sich für die Seite $s$ des regulären $n$-Eckes, wenn der Radius $c$ des umgeschriebenen Kreises gleich 1 angenommen wird,
\[
s = 2 \sin\frac{\pi}{n} = \sqrt{2\left(1 - \cos\frac{2\pi}{n}\right)}.
\]
Die übrigen Grössen $r$ stehen in derselben Beziehung zu den anderen Theilpunkten. Wegen dieser geometrischen Be-
%CONT>>
